\subsection{Selected components in the class-four quotient}
\label{subsec:selected-class-four}

Let $\Sigma$ be a set of even characters such that
$X^{[e]}$ is finite for each $e\in\Sigma$. In the
canonical quotient $Q_4$, write $B_4$ for the image
of $B=\ker(H\to W)$. Its commutator subgroup is
central, so $B_4$ has class at most two. Since $p$
is odd, its class-two BCH Lie algebra
$\mathfrak b_4=\log B_4$ is defined integrally.
Choose a tame $\Delta$ complement and put
\[
 J_\Sigma^{(4)}=
 \left\langle\!\left\langle
 \exp\!\left(\bigoplus_{\substack{e\text{ even}\\e\notin\Sigma}}
                         \mathfrak b_4^{[e]}\right)
 \right\rangle\!\right\rangle_{Q_4},\qquad
 Q_{4,\Sigma}=Q_4/J_\Sigma^{(4)},
\]
where the normal closure is closed. This definition
includes the omitted even central directions of
relative weight four.

\begin{proposition}\label{prop:selected-class-four}
The subgroup $J_\Sigma^{(4)}$ is independent of the
tame complement and is fully Galois-stable. The
canonical quotient $Q_{4,\Sigma}$ is detected by the
outer action, maps onto $Q_{3,\Sigma}$, and has
\[
             Q_{4,\Sigma}^{\rm ab}=(U/D_0)\oplus T_\Sigma.
\]
If $-2,-4\in\Sigma$, the quotient has nilpotency class
exactly four. In particular this holds for $\Sigma_{30}$,
and its detected field strictly contains the field
of $Q_{3,\Sigma_{30}}$.
\end{proposition}

\begin{proof}
For an arithmetic element $g$, the tame complements
$\Delta$ and $g\Delta g^{-1}$ in the same pro-$p$
extension are conjugate by some $h\in H$.
The element $h^{-1}g$ centralizes $\Delta$ pointwise:
it normalizes it and acts trivially on its image in
the abelian quotient $\operatorname{Gal}(F_\infty/\mathbf Q)$.
It therefore preserves every isotypic submodule of
$\mathfrak b_4$. Conjugation by $h$ preserves the
$Q_4$-normal closure of the displayed generators.
This proves full arithmetic invariance; the same
argument compares two tame complements. The image
in $Q_3$ of the even part of $\mathfrak b_4$ is exactly
$\bar E$, so the induced quotient maps onto the
previous $Q_{3,\Sigma}$. Detection follows from the
surjection $Q_4\twoheadrightarrow Q_{4,\Sigma}$.

The map $\mathfrak b_4\to Q_4^{\rm ab}$ is linear,
because its BCH correction is a commutator.
Its even image is exactly $T/Z_0$, and it commutes
with the integral $\Delta$ idempotents. Thus the
abelian image of $J_\Sigma^{(4)}$ is precisely the
omitted even summand of $T/Z_0$; taking its normal
closure adds no abelian directions. For a selected
character the single-component theorem gives $Z_e=0$,
and hence $(Z_0)_e=0$. The exact abelianization
$Q_4^{\rm ab}=(U/D_0)\oplus(T/Z_0)$ proves the assertion.

For strictness, take distinct positive odd weights
$m,n\geq3$ with $m\equiv n\equiv-1\pmod{p-1}$.
Theorem~\ref{thm:canonical-class-four} embeds $Q_4$
as its compact image in the full rational free Lie
target through generator length four. Choose
$\Delta$-eigen free generators in that target. Such
generators are obtained by applying the appropriate
idempotent to a lift of each linear generator; they
have unchanged linear terms. The resulting substitution
is a filtered automorphism, with its inverse constructed
successively in lengths two, three and four.

There is therefore an exact $\Delta$-equivariant Lie
homomorphism to the free class-four Lie algebra on
$x,y$, sending the generators of weights $m,n$ to
$x,y$ and every other generator to zero. Give $x,y$
the character $\omega^{-1}$. Exponentiation and the
actual target embedding give a group homomorphism
\[
       \Psi:Q_4\longrightarrow
           \exp\bigl(\operatorname{Lie}(x,y)/\gamma_5\bigr).
\]
The image of $B_4$ has zero linear term. Its even
Lie part consequently has only the characters $-2$
in length two and $-4$ in length four. Since these
characters are selected, all omitted even isotypes
map to zero, and $\Psi$ factors through $Q_{4,\Sigma}$.

Choose an integral $w\in W^{[-1]}$ whose moment
coordinates at $m,n$ are nonzero. A rational
$\Lambda$-basis vector, after clearing a bounded
denominator, supplies such a $w$. The linear images
of $w,\gamma w$ have determinant
\[
                   \lambda_m\lambda_n(u^n-u^m)\ne0.
\]
For actual lifts $a,b\in H$, denote these independent
linear images by $A,B$. The length-four image of
$[a,[a,[a,b]]]$ is $[A,[A,[A,B]]]$, which is nonzero
in the free Lie algebra on the independent pair
$A,B$. It survives in $Q_{4,\Sigma}$, proving class
four and strictness over its class-three quotient.
\end{proof}

\begin{corollary}[Class-four component detection]
\label{thm:class-four-components}
Let $p$ be odd and let $\Sigma$ be any set of even exponents
such that $X^{[e]}$ is finite for every $e\in\Sigma$.
The subgroup $J_\Sigma^{(4)}$ is independent of the tame
complement and stable under the full arithmetic Galois group.
The canonical field $L_\Sigma^{(4)}$ defined by $Q_{4,\Sigma}$
is Galois over $\mathbb Q$, and
\[
 L_\Sigma^{(2)}\subseteq L_\Sigma^{(3)}
   \subseteq L_\Sigma^{(4)}\subseteq\Omega_p^*,\qquad
 Q_{4,\Sigma}^{\rm ab}=(U/D_0)\oplus T_\Sigma.
\]
Here $L_\Sigma^{(3)}$ is the class-three field of
Corollary~\ref{thm:class-three-components}.
For $-2,-4\in\Sigma$, the quotient $Q_{4,\Sigma}$ has class
exactly four and $L_\Sigma^{(3)}\subsetneq L_\Sigma^{(4)}$.
\end{corollary}
This follows from Proposition~\ref{prop:selected-class-four}.
